\documentclass[a4paper,11pt]{article}
\usepackage[utf8]{inputenc}
\usepackage[T1]{fontenc}
\usepackage[french]{babel}
\usepackage{geometry}
\usepackage{xcolor}
\usepackage{tikz}

% Configuration des marges
\geometry{hmargin=2.5cm,vmargin=2.5cm}

\title{\textbf{Exercice de Cybersécurité}\\ \large ``Le vol des bijoux du Louvre''}
\author{}
\date{}

\begin{document}

\maketitle

\section{Objectif Pédagogique}
Sensibiliser les participants aux risques cyber liés :
\begin{itemize}
    \item à l’ingénierie sociale,
    \item au phishing ciblé (\textit{spear phishing}),
    \item à la compromission des accès internes,
    \item et à la gestion d'incident.
\end{itemize}

\section{Scénario}
[cite_start]Dans la nuit du \textbf{12 au 13 avril}, un vol spectaculaire a eu lieu au Louvre : la parure de diamants dite \textbf{``Les Larmes de Cléopâtre''} a disparu[cite: 9].

\begin{itemize}
    [cite_start]\item Aucune effraction physique n’a été constatée[cite: 10].
    [cite_start]\item Le système d’alarme n’a jamais été déclenché[cite: 10].
    [cite_start]\item Les accès numériques au coffre ont été validés avec les identifiants d’un cadre du département des antiquités[cite: 10].
\end{itemize}

La direction soupçonne une \textbf{attaque informatique} ayant permis de neutraliser la sécurité. [cite_start]Les équipes doivent maintenant reconstituer la chaîne d’attaque[cite: 11, 12].

\section{Éléments à analyser (Indices)}

\subsection{1. L'e-mail suspect (Phishing)}
\begin{quote}
    \textbf{Objet :} Mise à jour obligatoire – Partenariat UNESCO \\
    \textbf{Expéditeur :} julien.marcus@unesco-partnership.org \\
    
    Bonjour,\\
    Nous renouvelons notre partenariat et avons besoin de la liste complète des objets précieux enregistrés cette année au Louvre. Merci de compléter le document ci-joint.\\
    
    Cordialement,\\
    J. Marcus\\
    Département culturel – UNESCO\\
    
    [cite_start]\textit{[Pièce jointe : partenariat\_UNESCO.xlsm (Contient une macro malveillante)]} [cite: 19-24]
\end{quote}

\subsection{2. Extraits de logs réseau et systèmes}

[cite_start]\textbf{Log 1 – Proxy} [cite: 26]
\begin{itemize}
    \item Requête : \texttt{POST http://update-secure-cloud.com/upload.php}
    \item Utilisateur : \texttt{p.durand}
    \item Heure : 12/04 – 18:42
    [cite_start]\item Taille : 4,3 Mo [cite: 27-30]
\end{itemize}
    
[cite_start]\textbf{Log 2 – Système} [cite: 31]
\begin{itemize}
    \item Événement : Connexion réussie badge-numerique / \texttt{p.durand}
    \item Heure : 12/04 – 23:14
    [cite_start]\item Alerte : IP source interne inhabituelle \texttt{10.4.18.23} [cite: 32-34]
\end{itemize}

[cite_start]\textbf{Log 3 – Caméras} [cite: 35]
\begin{itemize}
    \item Événement : Perte de signal (salle des joyaux)
    \item Heure : 12/04 – 23:13 – 23:17
    [cite_start]\item Cause déclarée : mise à jour automatique [cite: 36-38]
\end{itemize}

\subsection{3. Organigramme du musée}
\begin{center}
\begin{tikzpicture}[
    node distance=1.5cm,
    every node/.style={rectangle, draw, rounded corners, align=center, minimum height=1cm, minimum width=2.5cm},
    level 1/.style={sibling distance=5cm},
    level 2/.style={sibling distance=3cm}
]

\node[fill=orange!20] {Directeur}
    child { node[fill=blue!20] {Accueil} }
    child { node[fill=red!20] {Antiquités}
        child { node[fill=red!20] {Département\\ d'art} }
    }
    child { node[fill=green!20] {Sécurité}
        child { node[fill=green!20] {Département\\ des salles} }
    };
\end{tikzpicture}
\end{center}

\section{Message interne – Équipe Sécurité}
[cite_start]\textit{Objet : Incident de sécurité – Analyse préliminaire et recommandations} [cite: 41]

Chers collaborateurs,

Dans la nuit du 12 au 13 avril, un incident majeur a été constaté. [cite_start]La parure « Les Larmes de Cléopâtre » a été dérobée sans déclenchement d’alarme [cite: 43-44].

[cite_start]\textbf{Premières constatations[cite: 47]:}
\begin{itemize}
    [cite_start]\item Une connexion interne a été réalisée à 23h14 depuis une adresse IP inhabituelle[cite: 48].
    [cite_start]\item Les caméras ont été désactivées pendant 4 minutes via une demande de ``mise à jour système''[cite: 49].
    [cite_start]\item Des traces d’exfiltration de données ont été identifiées sur un domaine externe[cite: 50].
\end{itemize}

[cite_start]\textbf{Consignes immédiates[cite: 51]:}
\begin{enumerate}
    [cite_start]\item Ne pas éteindre, redémarrer, ni modifier vos postes de travail[cite: 52].
    [cite_start]\item Signaler immédiatement toute activité inhabituelle[cite: 53].
    [cite_start]\item Ne pas communiquer publiquement sur l'incident[cite: 54].
\end{enumerate}

\section{Travail à faire (Résultat attendu)}
[cite_start]Vous devez fournir une \textbf{reconstitution complète de l’attaque} et proposer un \textbf{plan d’action concret} pour éviter un nouvel incident [cite: 63-64].

\end{document}